\documentclass[11pt,a4paper]{article}
\usepackage{amsmath,amssymb}
\usepackage{booktabs}
\usepackage{float}
\usepackage[hidelinks]{hyperref}
% generated by extract.py from docs/interest-audit.md; do not edit
\newcommand{\subjectCorpora}{6}
\newcommand{\subjectKept}{952 to 1904}
\newcommand{\subjectRZero}{19 to 50}
\newcommand{\subjectROne}{1 to 2}
\newcommand{\subjectRTwo}{2 to 6}
\newcommand{\subjectClockStops}{0}
\newcommand{\subjectBaseRate}{2.1\% to 3.0\%}
\newcommand{\subjectTopKShare}{0.0\% to 6.0\%}
\newcommand{\subjectImpShare}{4.0\% to 9.5\%}
\newcommand{\subjectRhoU}{$-$0.03 to 0.29}
\newcommand{\subjectRhoImp}{$-$0.50 to $-$0.25}
\newcommand{\subjectRhoFoot}{0.56 to 0.80}
\newcommand{\subjectRhoUsers}{$-$0.57 to $-$0.19}
\newcommand{\subjectRandLo}{1.7\% to 3.3\%}
\newcommand{\subjectRandHi}{11.7\% to 15.0\%}
\newcommand{\subjectLoopI}{1.7\% to 8.3\%}
\newcommand{\subjectLoopFilt}{5.3\% to 23.1\%}
\newcommand{\subjectLoopSize}{0.0\% to 3.3\%}
\newcommand{\subjectLoopIAbove}{0}
\newcommand{\subjectLoopIBelow}{1}
\newcommand{\subjectLoopIWithin}{5}
\newcommand{\subjectLoopFiltAbove}{1}
\newcommand{\subjectLoopFiltBelow}{0}
\newcommand{\subjectLoopFiltWithin}{5}
\newcommand{\subjectRuns}{6}
\newcommand{\subjectFiltPromoted}{10 to 19}
\newcommand{\controlCorpora}{6}
\newcommand{\controlKept}{1045 to 1804}
\newcommand{\controlRZero}{7 to 14}
\newcommand{\controlROne}{0 to 1}
\newcommand{\controlRTwo}{0}
\newcommand{\controlClockStops}{0}
\newcommand{\controlBaseRate}{0.4\% to 1.3\%}
\newcommand{\controlTopKShare}{0.0\%}
\newcommand{\controlImpShare}{0.0\%}
\newcommand{\controlRhoU}{$-$0.29 to 0.09}
\newcommand{\controlRhoImp}{$-$0.69 to $-$0.60}
\newcommand{\controlRhoFoot}{0.77 to 0.91}
\newcommand{\controlRhoUsers}{$-$0.70 to $-$0.37}
\newcommand{\controlRandLo}{0.0\%}
\newcommand{\controlRandHi}{3.3\% to 6.7\%}
\newcommand{\controlLoopI}{0.0\%}
\newcommand{\controlLoopFilt}{0.0\% to 20.0\%}
\newcommand{\controlLoopSize}{0.0\%}
\newcommand{\controlLoopIAbove}{0}
\newcommand{\controlLoopIBelow}{0}
\newcommand{\controlLoopIWithin}{6}
\newcommand{\controlLoopFiltAbove}{4}
\newcommand{\controlLoopFiltBelow}{0}
\newcommand{\controlLoopFiltWithin}{2}
\newcommand{\controlRuns}{6}
\newcommand{\controlFiltPromoted}{10}
\newcommand{\hTwoTrivRecallOne}{0\%}
\newcommand{\hTwoTrivRecallTwo}{100\%}
\newcommand{\hTwoTrivFPR}{0.0\%}
\newcommand{\hTwoTrivJ}{0.500 to 0.857}
\newcommand{\hTwoRestRecallOne}{100\%}
\newcommand{\hTwoRestRecallTwo}{0\% to 100\%}
\newcommand{\hTwoRestFPR}{99.2\% to 99.7\%}
\newcommand{\hTwoRestJ}{$-$0.494 to 0.006}
\newcommand{\hTwoCombRecallOne}{100\%}
\newcommand{\hTwoCombRecallTwo}{100\%}
\newcommand{\hTwoCombFPR}{99.2\% to 99.7\%}
\newcommand{\hTwoCombJ}{0.003 to 0.008}
\newcommand{\hOneVariants}{12}
\newcommand{\hOneRankings}{15}
\newcommand{\hOneEnrichedAny}{0}
\newcommand{\reportSha}{652d4a920ece15d2}
\newcommand{\preregSha}{9aadc02723a75b65}
\newcommand{\branchHead}{0324ff7}
\newcommand{\dupIdentical}{byte-identical}
\newcommand{\verdictRows}{7}
\newcommand{\nearDupK}{50}
\newcommand{\nearDupDraws}{200}
\newcommand{\subjectNearDupIZero}{0\% to 16\%}
\newcommand{\subjectNearDupIZeroMedian}{12\%}
\newcommand{\subjectNearDupUZero}{0\% to 16\%}
\newcommand{\subjectNearDupUZeroMedian}{2\%}
\newcommand{\subjectNearDupImp}{0\% to 12\%}
\newcommand{\subjectNearDupImpMedian}{0\%}
\newcommand{\subjectNearDupRand}{1\% to 3\%}
\newcommand{\subjectNearDupRandMedian}{3\%}
\newcommand{\subjectNearDupAboveRandom}{5}
\newcommand{\controlNearDupIZero}{0\% to 36\%}
\newcommand{\controlNearDupIZeroMedian}{14\%}
\newcommand{\controlNearDupUZero}{0\% to 16\%}
\newcommand{\controlNearDupUZeroMedian}{2\%}
\newcommand{\controlNearDupImp}{0\% to 12\%}
\newcommand{\controlNearDupImpMedian}{4\%}
\newcommand{\controlNearDupRand}{2\% to 3\%}
\newcommand{\controlNearDupRandMedian}{3\%}
\newcommand{\controlNearDupAboveRandom}{5}


\newcommand{\Iz}{I_0}
\newcommand{\Uz}{U_0}
\newcommand{\arxiv}[1]{\href{https://arxiv.org/abs/#1}{arXiv:#1}}
\newcommand{\repo}[2]{\href{https://github.com/carlok/euclean/#1/main/#2}{\texttt{#2}}}

\title{Interestingness against a theory with known restatements\\[2pt]
\large A pre-registered audit of Patel et al.\ (2026) on an anonymized formal theory}
\author{Carlo Perassi}
\date{November 2026}

\begin{document}
\maketitle

\begin{abstract}
\noindent
Patel et al.~\cite{patel} propose scoring a theorem by how long its proof is relative to how
short its statement is, and use the score to decide which machine-generated
theorems enter a growing formal library. A known failure of compression- and
frequency-based scores is that they reward \emph{restatements}: results that
look new but only carry one axiom's conclusion forward. We test the metric on a
small anonymized first-order theory where every restatement can be labelled
exactly from proof dependencies, with a second theory as a control, under a
pre-registration frozen before any score was computed and by an executor with
no access to the theory's interpretation. The metric does not reward
restatements; if anything it avoids them. It also does not track usefulness on
this material: it follows how many axioms a result needs, and runs against how
often the result is reused. The paper's own filters either miss the
restatements or catch them only by rejecting nearly everything. An open
implementation of the metric and the labelled corpora are released with this
note.
\end{abstract}

\section{Setting}

The test bed is a research pipeline that derives consequences of a first-order
theory, checks every derivation with the Lean~4 kernel, and then tries to
recover structure from the resulting corpus without being told what the
symbols mean. The theory is given as one sort, two relation symbols of fixed
arity and ten axioms, all under opaque names. Its interpretation exists but is
held outside the analysed tree and is consulted only after results are frozen.
A second, richer theory under the same treatment serves as a control.

Each corpus is a set of kernel-verified statements with their proof terms and
the full citation graph among them. Because proofs are built as terms and
only checked by Lean, the dependency graph is exact: for every statement we know
which axioms it ultimately rests on.

An earlier stage of this project found that its own scores for invented
concepts, based on compression and frequency, favoured a pattern that turned
out to be one axiom's conclusion recurring across hundreds of proofs. Surface
similarity would not have exposed it; the axiom-dependency footprint did.

\section{What is already known}

Interestingness as a search criterion goes back to Lenat's AM~\cite{lenat} and
to Colton's HR~\cite{colton}, which ranked conjectures and concepts with
hand-designed measures of novelty, surprise and applicability. Recent work
returns to the question at the scale of Lean and Mathlib.
LeanConjecturer~\cite{leanconjecturer} generates conjectures for theorem
proving; Porto~\cite{porto} asks whether a verified theorem is actually new,
combining library search, triviality tactics and how far apart two proofs' premise sets are.

Patel et al.~\cite{patel} define \emph{intrinsic interestingness} as
\[
  I(T\mid P) = 100\,\frac{V(T\mid P)}{L(T\mid P)},
\]
where $V$ is the proof cost of $T$ given premises $P$, counted in source LoC,
and $L$ is the size in characters of the statement plus the definitions
it needs that $P$ does not already supply. With no premises this gives an
absolute score $\Iz(T)$. They report Spearman $\rho = 0.756$ against a utility
$\Uz(T) = |D(T)|\cdot V(T\mid\varnothing)$, where $D(T)$ is the set of results
citing $T$ directly, on Mathlib. The score is used as a reward for a conjecturing
model and as the promotion rule in a self-expanding library, with a triviality
filter (\texttt{assumption}, \texttt{rfl}, \texttt{simp}, \texttt{tauto},
\texttt{simp\_all}), a token-similarity restatement filter, and a language-model
de-duplication step.

\section{The question}

On Mathlib nobody can say exactly which theorems are restatements. Here we
can. The question is whether $\Iz$ ranks restatements highly, whether the
paper's filters catch them, whether the $\Iz$--$\Uz$ relation carries over, and
what a promotion loop driven by $\Iz$ ends up keeping.

Restatements are labelled from proof dependencies alone, before any score is
seen: R0, the statement is an axiom's conclusion up to canonical form; R1, its
dependency footprint is a single axiom and it follows by instantiation and
argument permutation; R2, every proof supporting it cites the axiom whose
conclusion it is. Everything else is N.

Five hypotheses were fixed in advance, each with a decision rule:
H1, top-50 enrichment of restatements under $\Iz$; H2, recall of the paper's
filters on R1 and R2, with their false-positive rate on N; H3, rank
correlations of $\Iz$ with $\Uz$, with this project's structural importance
ranking and with footprint size; H4, a six-round promotion loop with no
language model, candidates drawn from the corpus; H5, an exploratory corrected
score that decides nothing.

\section{What was found}

\paragraph{$\Iz$ does not reward restatements.}
Across all \hOneVariants{} variants of the metric (three ways of measuring $V$, two
ways of expanding premises, raw or canonicalised statements) not one ranking
puts more restatements in its top 50 than the corpus base rate. In the control
theory the top 50 contains none at all; in the subject, one permutation in three
is significantly depleted. The same holds for $\Uz$ and for this project's own
importance ranking: none of the \hOneRankings{} rankings tested is enriched.
The restatement trap that motivated the audit is not where this metric fails.

\paragraph{The paper's filters do not separate restatements.}
The triviality filter catches every propagated restatement (R2) and no
single-axiom one (R1), while flagging almost nothing else. The token-similarity
restatement filter, which the paper describes without a threshold, flags
\hTwoRestFPR{} of ordinary statements at the pre-registered stand-in: in a
small signature most statements share most tokens with some premise. Combined,
the two catch every restatement only by rejecting nearly the whole corpus, so
the pre-registered recall rule is met with no separation at all
(Youden's $J$ \hTwoCombJ). The control theory produced no R1 or R2 statements,
so H2 is undefined there.

\paragraph{$\Iz$ does not track usefulness here.}
The paper's headline relation does not carry over: $\rho(\Iz,\Uz)$ ranges over
\subjectRhoU{} in the subject and \controlRhoU{} in the control, against 0.756
on Mathlib. What $\Iz$ does follow is the size of a statement's axiom footprint
($\rho$ \subjectRhoFoot{}, control \controlRhoFoot). It runs \emph{against}
structural importance ($\rho$ \subjectRhoImp{}, control \controlRhoImp) and
against direct reuse ($\rho$ \subjectRhoUsers{}, control \controlRhoUsers). On
this material the score favours results that are expensive to derive and
little used.

\paragraph{Promotion by $\Iz$ is indistinguishable from random.}
In a six-round loop promoting ten statements per round, promotion by $\Iz$
never puts more restatements among the promoted than random promotion does: it
falls inside the random band in \subjectLoopIWithin{} of \subjectRuns{} subject
runs and below it in \subjectLoopIBelow, and inside it in \controlLoopIWithin{} of
\controlRuns{} control runs. Adding the paper's restatement filter to the loop
moves the share \emph{up}: above the random band in \subjectLoopFiltAbove{} of
\subjectRuns{} subject runs and \controlLoopFiltAbove{} of \controlRuns{} control
runs. Removing almost every candidate leaves a residue richer in restatements.
That policy promoted only \subjectFiltPromoted{} statements per run in the
subject and \controlFiltPromoted{} in the control, so this is a warning, not a
measurement.

\paragraph{Its top ranks hold near-duplicates (post hoc).}
This was not pre-registered. Reading the top-ranked statements after the audit
was frozen showed several that differ from another top-ranked statement only in
one hypothesis. Measured structurally, as the share of the top \nearDupK{} that
has another top-\nearDupK{} entry one hypothesis apart, it is a median of
\subjectNearDupIZeroMedian{} under $\Iz$ in the subject and
\controlNearDupIZeroMedian{} in the control (per corpus, \controlNearDupIZero),
against \subjectNearDupRandMedian{} for random draws; $\Iz$ exceeds random in
\subjectNearDupAboveRandom{} of 6 subject and \controlNearDupAboveRandom{} of 6
control corpora. $\Uz$ and structural importance show it less: medians of
\subjectNearDupUZeroMedian{} and \controlNearDupUZeroMedian{} for $\Uz$,
\subjectNearDupImpMedian{} and \controlNearDupImpMedian{} for importance. The
paper's language-model filter, which removes equivalent conjectures and
enforces diversity across theorem families, is aimed at this, and the audit did
not reproduce it: the effect is evidence that the filter does real work, not
that it is missing.

\paragraph{Method.}
Three choices did the work. Labels were computed from exact proof dependencies
and frozen by hash before any score existed, so no verdict could influence them. The
executor never read the interpretation; it did glimpse two wordlist entries in
guard output, which the pre-registration discloses. Wording that touched the
anonymization was corrected by a separate session that could see it. A control
theory separated what belongs to the metric from what belongs to one theory:
several effects appear in both, which is what makes them claims about the
metric.

\section{Limits}

\begin{itemize}
\item Restatements are rare, \subjectBaseRate{} of the subject corpus and
  \controlBaseRate{} of the control, so a top-50 list holds one or two by
  chance. H1 can detect strong enrichment and depletion; it cannot detect a
  mild preference.
\item $V$ counts LoC of emitted proof terms under one fixed layout, plus two
  alternatives (term size, a fixed tactic ladder). The paper counts tactic
  proofs written by people and models. All three give the same H1 verdict,
  which bounds the style objection but does not remove it; H3 is reported for
  the primary variant only.
\item Statements are short and identifiers anonymized, so only ranks within a
  corpus are compared, never absolute values with the paper's.
\item The language-model de-duplication step was not reproduced.
\item Two theories, six corpora each. The conclusions are about this metric on
  small first-order theories, not about Mathlib.
\item The restatement filter's threshold is not given in the paper; the
  pre-registered stand-in (token-set Jaccard 0.70, the value the paper uses
  for a different check) was fixed in advance and a sweep is reported, but a
  different threshold is a different filter.
\end{itemize}

\section{Code and data}

Everything is in the public repository
\begin{center}\url{https://github.com/carlok/euclean}\end{center}
\begin{itemize}
\item \repo{blob}{docs/interest-audit.md}: the audit report.
\item \repo{blob}{docs/interest-prereg.md}: its pre-registration.
\item \repo{tree}{kit/}, an implementation of the paper's quantities over a plain JSON
  library format, with the twelve corpora and their frozen labels. Every choice
  the paper leaves open is a named parameter there. Its tests recompute the
  audit's values from the exported corpora alone and require exact agreement.
  The near-duplicate measurement above is one of its examples.
\item \repo{tree}{paper/interest-audit/}, this note: its source, and the script that
  reads every number in it from the frozen report and from that example's
  output.
\end{itemize}

\section*{Acknowledgements}

The pipeline, the audit and this note were produced with AI coding agents
(Claude, by Anthropic) working under the author's direction. The author is
responsible for the content.

\clearpage
\appendix
\section{Numbers}

All values are ranges over three identifier permutations and two base seeds,
read by script from the frozen report, itself generated from run artifacts.

\smallskip
{\small\begin{tabular}{@{}ll@{}}
report sha256 prefix & \texttt{\reportSha} \\
pre-registration sha256 prefix & \texttt{\preregSha} \\
audit head commit & \texttt{\branchHead} \\
\end{tabular}}

\begin{table}[H]
\centering\footnotesize
\begin{tabular}{lll}
\toprule
 & subject & control \\
\midrule
corpora & \subjectCorpora & \controlCorpora \\
statements per corpus & \subjectKept & \controlKept \\
R0 per corpus & \subjectRZero & \controlRZero \\
R1 per corpus & \subjectROne & \controlROne \\
R2 per corpus & \subjectRTwo & \controlRTwo \\
restatement base rate & \subjectBaseRate & \controlBaseRate \\
corpora with a clock stop & \subjectClockStops & \controlClockStops \\
\midrule
top-50 restatement share, $\Iz$ & \subjectTopKShare & \controlTopKShare \\
top-50 restatement share, importance & \subjectImpShare & \controlImpShare \\
\midrule
$\rho(\Iz, \Uz)$, $\Uz>0$ & \subjectRhoU & \controlRhoU \\
$\rho(\Iz, \text{importance})$ & \subjectRhoImp & \controlRhoImp \\
$\rho(\Iz, \text{footprint size})$ & \subjectRhoFoot & \controlRhoFoot \\
$\rho(\Iz, \text{direct users})$ & \subjectRhoUsers & \controlRhoUsers \\
\midrule
loop, random band (low edge) & \subjectRandLo & \controlRandLo \\
loop, random band (high edge) & \subjectRandHi & \controlRandHi \\
loop, promoted by $\Iz$ & \subjectLoopI & \controlLoopI \\
loop, $\Iz$ with restatement filter & \subjectLoopFilt & \controlLoopFilt \\
loop, promoted by proof size & \subjectLoopSize & \controlLoopSize \\
\midrule
near-duplicate share, top \nearDupK, $\Iz$ & \subjectNearDupIZero & \controlNearDupIZero \\
near-duplicate share, top \nearDupK, $\Uz$ & \subjectNearDupUZero & \controlNearDupUZero \\
near-duplicate share, top \nearDupK, importance & \subjectNearDupImp & \controlNearDupImp \\
near-duplicate share, random \nearDupK & \subjectNearDupRand & \controlNearDupRand \\
\bottomrule
\end{tabular}
\caption{Corpora, H1, H3, H4, and the post-hoc near-duplicate measurement.}
\end{table}

\begin{table}[H]
\centering\small
\begin{tabular}{lllll}
\toprule
filter (subject) & recall R1 & recall R2 & FPR on N & Youden's $J$ \\
\midrule
triviality & \hTwoTrivRecallOne & \hTwoTrivRecallTwo & \hTwoTrivFPR & \hTwoTrivJ \\
restatement & \hTwoRestRecallOne & \hTwoRestRecallTwo & \hTwoRestFPR & \hTwoRestJ \\
combined & \hTwoCombRecallOne & \hTwoCombRecallTwo & \hTwoCombFPR & \hTwoCombJ \\
\bottomrule
\end{tabular}
\caption{H2 on the subject theory. The control has no R1 or R2 statements.}
\end{table}

A duplicate build of one control corpus, run under different machine load,
was \dupIdentical{} to the original; no corpus in the audit had a generation
cut short by the engine's time budget.

\begin{thebibliography}{9}

\bibitem{patel} N.~Patel, A.~Rammal, A.~Hayat, R.~Munos, J.~Kempe.
\emph{Learning to Discover Interesting Mathematics}. \arxiv{2609.28603}, 2026.

\bibitem{lenat} D.~Lenat. \emph{AM: An Artificial Intelligence Approach to
Discovery in Mathematics as Heuristic Search}. PhD thesis, Stanford University,
1976.

\bibitem{colton} S.~Colton. \emph{Automated Theory Formation in Pure
Mathematics}. Springer, 2002.

\bibitem{leanconjecturer} N.~Onda, K.~Kasaura, Y.~Oriike, M.~Taniguchi,
A.~Sannai, S.~Sonoda. \emph{LeanConjecturer: Automatic Generation of
Mathematical Conjectures for Theorem Proving}. \arxiv{2506.22005}, 2025.

\bibitem{porto} A.~Porto. \emph{Beyond Correctness: Toward Automated Novelty
Verification with Lean 4}. \arxiv{2608.14669}, 2026.

\end{thebibliography}

\end{document}
